% ==================== CONCLUSION ====================
% STYLE GUIDELINES:
% - Maximum 2 pages (550 words maximum)
% - Summarize work carried out and main contributions
% - Highlight difficulties encountered
% - Propose prospects for future development
% - Font: Times New Roman 12pt, line spacing 1.5
% - Section titles: left-aligned, bold (if used)
% ====================================================================
\phantomsection
{\huge\bfseries General Conclusion}
\addcontentsline{toc}{chapter}{General Conclusion}
\markboth{General Conclusion}{}     % <-- add this

\vspace{1cm}

\initial{A}s part of this \textbf{four-month} internship at Novation City,
IoT-REAP was designed and implemented, a secure industrial remote operations and training
platform for Industry~4.0 learning. The platform unifies six functional domains-session
management, VM provisioning, remote desktop access, camera supervision, IoT device control,
and structured training management-within a role-differentiated interface serving Engineers,
Teachers, and Administrators. The objective was to eliminate physical laboratory dependency
by enabling users to launch sessions, interact with real equipment, follow guided learning
paths, and monitor activity from a single web interface.

The development approach was iterative and architecture-driven, following Agile-inspired
delivery cycles. The backend leverages Laravel organized around FormRequest validation,
thin controllers, service-layer logic, and repository-based data access. The frontend uses
React with TypeScript for typed API integrations and reusable components. Role-based access
ensures each profile encounters only contextually relevant actions, forming a solid foundation
for external system integration~{\footnotesize\cite{laravel-docs,react-docs,typescript-docs}}.

A major contribution was the remote lab session lifecycle. Proxmox integration enables
on-demand VM provisioning and lifecycle control, while Guacamole-based browser access
supports RDP, VNC, and SSH without local client installation. Tokenized access patterns,
expiration-aware handling, and resource cleanup prevent orphaned VMs and maintain stability,
bridging cloud orchestration and industrial training~{\footnotesize\cite{proxmox-ve-api,guacamole-manual}}.

A second contribution was integrating real-time hardware context into training sessions.
Camera services and gateway integrations enable remote equipment observation, while PTZ
control routes through structured, auditable MQTT command channels. USB device reservation
workflows attach specific hardware to live sessions in a controlled manner, positioning
IoT-REAP as an operational remote laboratory rather than a conventional e-learning
portal~{\footnotesize\cite{mqtt-oasis,usbip-docs}}.

The platform supports instructor-driven training path management. Teachers define learning
sequences linked to specific VM sessions, synchronizing pedagogical structure with practical
execution. Learners progress through guided material while interacting with realistic
industrial environments, developing operational competencies through direct engagement.

Delivering these capabilities required navigating significant engineering challenges.
Coordinating heterogeneous technologies-Proxmox's asynchronous operations, state-sensitive
remote desktop tunnels, MQTT topic discipline, and variable camera streaming-required
isolating external provider logic in dedicated service classes, standardizing integration
contracts, and improving observability through explicit logging. Preserving architectural
quality amid rapid feature expansion demanded disciplined coding standards and incremental
API documentation to prevent business logic leakage across layers. Balancing delivery velocity
with security requirements necessitated explicit authorization boundaries, input validation,
role scoping, and regression tests on security-critical paths, producing a higher-assurance,
auditable workflow.

Future extensions include completing high-impact API integrations and advancing frontend
parity for administrative workflows. Deeper analytics on session usage, infrastructure
efficiency, and learner progression would support data-driven decisions. Predictive scheduling
and optimized resource allocation could reduce wait times, while enhanced zero-trust controls,
compliance-oriented audit reporting, and mobile-optimized interfaces would broaden operational
reach.

This internship produced tangible technical outcomes while developing competencies across
full-stack architecture, distributed system integration, and security-aware design. IoT-REAP
delivers a technically sound foundation connecting virtual infrastructure, remote hardware
interaction, and structured learning. Core requirements were verified through unit and feature
tests, confirming specification alignment. The system demonstrates the feasibility of secure,
remote Industry~4.0 training at scale, with an architecture designed to accommodate a growing
operational roadmap.

\cleardoublepage